% problem-000352 6 木犀草素溶解热力学
\section*{6. 木犀草素溶解热力学}
（图：）

\subsection*{6-1}
实验测得， $4 0 ~ ^ { \circ } \mathrm { C }$ 时⽊犀草素在⼄醇(EtOH)中的溶解度（以物质的量分数计）为 $2 . 6 8 \times 1 0 ^ { - 3 }$ 。已知$2 0 { \sim } 6 0 ~ ^ { \circ } \mathrm { C }$ ⽊犀草素在⼄醇中的溶解焓为20.28kJ $\mathrm { m o l ^ { - 1 } s }$ 。请估算 $2 5 ^ { \circ } \mathrm { C }$ 时⽊犀草素在⼄醇中的溶解度。

\subsection*{6-2}
$4 0 ~ ^ { \circ } \mathrm { C }$ 时，芦丁和⽊犀草素在⼄醇-⽔混合溶剂中的溶解度曲线如图6.1。根据图⽰，设计提纯⽊犀草素的具体⽅案。

\subsection*{6-3}
⽊犀草素有α和β两种晶相，在热⼒学上，β相⽐α相稳定。在恒压下，液态、α相和β相的摩尔Gibbs⾃由能 $G _ { \mathrm { m } }$ 随温度�的变化关系如图6.2所⽰。

\subsection*{6-3}
-1 请回答哪条线对应α相、β相，说明理由。

\subsection*{6-3}
-2 在相同温度下，α相和β相哪个熵更⼤？简述理由。

（图：）
图 6.1

（图：）
图 6.2
